% Source: https://artofproblemsolving.com/community/c3h2319138_200th_post_celebration_geometry
% Stable identity: aops-topic-2319138-q02
$\boxed{\text{False}}$. Adjacent angles share a vertex and a side, and have no interior points in common. Take rays $OX$, $OY$, $OZ$ with $\angle XOY=\angle YOZ=30^\circ$. The angles $\angle XOY$ and $\angle YOZ$ are adjacent and congruent, but together they make only $60^\circ$, so $OX$ and $OZ$ are not opposite rays and the angles do not form a linear pair.

\begin{asy}
usepackage("amsmath");
size(5cm);
defaultpen(fontsize(10pt));
pair O=(0,0);
draw(O--3*dir(0),linewidth(0.9),Arrow(5)); draw(O--3*dir(30),linewidth(0.9),Arrow(5)); draw(O--3*dir(60),linewidth(0.9),Arrow(5));
draw((-2.2,0)--O,dashed+gray(0.5));
draw(arc(O,0.9,0,30),blue); draw(arc(O,1.1,30,60),red);
label("$30^\circ$",1.25*dir(15),blue); label("$30^\circ$",1.45*dir(45),red);
label("$X$",3*dir(0),S); label("$Y$",3*dir(30),SE); label("$Z$",3*dir(60),NW); label("$O$",O,SW);
dot(O);
\end{asy}

(Adjacent congruent angles form a linear pair only when each of them is $90^\circ$.)
